\documentclass[11pt,a4paper]{article}
\usepackage[margin=1in]{geometry}
\usepackage[T1]{fontenc}
\usepackage[utf8]{inputenc}
\usepackage{lmodern}
\usepackage{microtype}
\usepackage{xcolor}
\usepackage{graphicx}
\usepackage{hyperref}
\usepackage{booktabs}
\usepackage{longtable}
\usepackage{array}
\usepackage{listings}
\usepackage{float}
\usepackage{fancyhdr}
\usepackage{titlesec}
\usepackage{setspace}
\usepackage{amsmath}
\usepackage{amssymb}
\usepackage{amsthm}

\hypersetup{
  colorlinks=true,
  linkcolor=blue!60!black,
  urlcolor=blue!60!black,
  citecolor=blue!60!black,
  linktoc=all,
  pdftitle={CS6160 Assignment 1 -- Problem 4: ZKP of Polynomial Evaluation},
  pdfauthor={Divyansh Sevta, Nishanth D.}
}
\lstdefinestyle{pythonstyle}{
  language=Python, basicstyle=\ttfamily\small,
  frame=single, breaklines=true,
  keywordstyle=\color{blue!70!black},
  commentstyle=\color{green!50!black},
  columns=fullflexible, showstringspaces=false
}
\lstdefinestyle{bashstyle}{
  language=bash, basicstyle=\ttfamily\small,
  frame=single, breaklines=true,
  columns=fullflexible, showstringspaces=false
}
\lstdefinestyle{jsonstyle}{
  basicstyle=\ttfamily\small,
  frame=single, breaklines=true,
  columns=fullflexible, showstringspaces=false
}
\pagestyle{fancy}
\fancyhf{}
\lhead{CS6160: Advanced Topics in Cryptology}
\rhead{Assignment 1 -- Problem 4}
\cfoot{\thepage}
\titleformat{\section}{\large\bfseries}{\thesection.}{0.5em}{}
\titleformat{\subsection}{\normalsize\bfseries}{\thesubsection}{0.5em}{}
\setstretch{1.15}
\setcounter{tocdepth}{2}

\newtheorem{theorem}{Theorem}
\newtheorem{lemma}{Lemma}
\newtheorem{definition}{Definition}

\begin{document}

\begin{titlepage}
\centering
\vspace*{0.6cm}
{\Large\bfseries Indian Institute of Technology Hyderabad\par}
\vspace{0.35cm}
{\large Department of Computer Science and Engineering\par}
\vspace{0.2cm}
{\large\textit{CS6160: Advanced Topics in Cryptology}\par}
\vspace{0.7cm}
\rule{0.78\textwidth}{0.6pt}\par
\vspace{0.6cm}
{\LARGE\bfseries Programming Assignment 1 --- Problem 4\par}
\vspace{0.3cm}
{\Large Zero-Knowledge Proof of Knowledge and\\ Correct Evaluation of a Secret Polynomial\par}
\vspace{0.6cm}
\rule{0.78\textwidth}{0.6pt}\par
\vspace{0.9cm}
\begin{minipage}{0.72\textwidth}
\centering
\begin{tabular}{@{}rl@{}}
\textbf{Group Member 1:} & Divyansh Sevta \\
\textbf{Roll Number:}    & CS25MTECH14018 \\
\textbf{Group Member 2:} & Nishanth D. \\
\textbf{Roll Number:}    & CS25MTECH14020 \\
\end{tabular}
\end{minipage}
\vspace{0.8cm}

\vspace{0.5cm}
{\textbf{Instructor:} Dr.~Maria Francis \qquad \textbf{Semester:} Spring 2026}
\end{titlepage}

\tableofcontents
\newpage

\begin{abstract}
We design and implement a zero-knowledge proof system by which a prover
convinces a verifier that (i) the prover knows the coefficients
$a_0,\dots,a_d \in \mathbb{Z}_q$ of a secret polynomial $P(x)$ that match
public Pedersen commitments $C_0,\dots,C_d$, and (ii) $P(z) \equiv y \pmod q$
for a publicly supplied point $z$ and claimed value $y$. The protocol is
a parallel per-coefficient $\Sigma$-protocol with a single shared
challenge, augmented by a scalar evaluation-consistency check that binds
the opening of each commitment to the claimed evaluation. We prove
completeness, $2$-special-soundness with a full-witness extractor, and
perfect honest-verifier zero-knowledge. Both a multi-round interactive
protocol and a Fiat--Shamir non-interactive variant are implemented in
pure Python 3 with only \texttt{hashlib}, \texttt{secrets}, and standard
built-ins; no external cryptographic library is used.
\end{abstract}

\section{Introduction}

Polynomials are a fundamental object in modern cryptography: polynomial
commitments (KZG, Bulletproofs, FRI) underpin SNARKs, STARKs, and
verifiable secret-sharing. In this assignment we consider the simplest
instance of this pattern. The prover holds a secret polynomial of degree
$d$ over $\mathbb{Z}_q$; the public parameters include Pedersen
commitments to each coefficient and a challenge point $z$. The prover
must convince a verifier that the polynomial evaluates to a claimed
value $y$ at $z$, without revealing any coefficient.

Our construction is deliberately elementary: we do not rely on pairings
or on polynomial interpolation tricks. Instead we run $d+1$ parallel
Schnorr proofs of knowledge of the individual Pedersen openings, share a
single challenge across them, and attach one extra scalar equation that
binds the coefficients to $y$. This yields an efficient $\Sigma$-protocol
that provably meets the assignment's soundness requirement: from two
accepting transcripts with distinct challenges, one can extract the
\emph{full} witness $(a_0,\dots,a_d, r_0,\dots,r_d)$.

\section{Preliminaries}

\subsection{Cyclic group and discrete log assumption}

Let $p$ be a prime and $q$ a prime divisor of $p-1$. Let $G_q \subseteq
\mathbb{Z}_p^{*}$ be the unique subgroup of order $q$ with public
generator $g$. We additionally fix an \emph{independent} generator
$h \in G_q$ such that $\log_g h$ is unknown. All exponent arithmetic is
in $\mathbb{Z}_q$; all group-element arithmetic is in $\mathbb{Z}_p^{*}$.

\begin{definition}[Discrete Logarithm Problem]
Given $g, y \in G_q$, find $x \in \mathbb{Z}_q$ with $g^x = y \pmod p$.
\end{definition}

We assume DLP is hard in $G_q$. The binding property of the Pedersen
commitment below reduces to this assumption.

\subsection{Pedersen commitment}

\begin{definition}[Pedersen Commitment]
For $v \in \mathbb{Z}_q$ and $r \xleftarrow{\$} \mathbb{Z}_q$,
\[
\mathrm{Com}(v;\,r) \;=\; g^{v}\,h^{r} \pmod p.
\]
\end{definition}

\paragraph{Perfect hiding.} The map $r \mapsto h^{r}$ is a bijection from
$\mathbb{Z}_q$ onto $G_q$. For any fixed $v$, the distribution of
$g^{v}h^{r}$ over uniform $r$ is the uniform distribution on $G_q$,
independent of $v$. Hence Pedersen is \emph{perfectly} hiding.

\paragraph{Computational binding.} Suppose an adversary opens a single
commitment $C$ two ways: $C = g^{a}h^{r} = g^{a'}h^{r'}$ with
$(a,r)\neq(a',r')$. Then $g^{\,a-a'} = h^{\,r'-r}$, giving
$\log_g h = (a-a')(r'-r)^{-1} \bmod q$. Under the DL assumption this is
infeasible, so Pedersen commitments are computationally binding.

\section{Problem Statement}

\paragraph{Public inputs.} $(G_q, q, g, h)$; polynomial degree $d$;
coefficient commitments $C_0,\dots,C_d$; evaluation point
$z \in \mathbb{Z}_q$; claimed value $y \in \mathbb{Z}_q$; rounds $k$.

\paragraph{Witness.} $\mathbf{w} = (a_0,\dots,a_d,\; r_0,\dots,r_d)
\in \mathbb{Z}_q^{2(d+1)}$.

\paragraph{Relation.}
\[
\mathcal{R} \;=\; \Bigl\{\bigl((C_0,\dots,C_d,z,y);\;\mathbf{w}\bigr)
\;:\; C_i = g^{a_i}h^{r_i}\ \forall i,\
\textstyle\sum_{i=0}^{d} a_i z^{i} \equiv y \pmod q \Bigr\}.
\]

A valid protocol must satisfy completeness, computational soundness, and
(honest-verifier) zero-knowledge. A knowledge-sound protocol, moreover,
must admit a \emph{knowledge extractor} that, given sufficient
transcripts, outputs a witness in $\mathcal{R}$.

\section{Why Naive Aggregation Fails}

A tempting shortcut is to aggregate commitments as
$C_{\mathrm{eval}} = \prod_i C_i^{z^i}$, observe that
$C_{\mathrm{eval}} = g^{P(z)} h^{R}$ with $R = \sum_i r_i z^{i}$, and
run a single Schnorr PoK of $R$ over $D = C_{\mathrm{eval}}\cdot g^{-y}$.

The extractor for this single-aggregate protocol outputs only the scalar
$\hat{R} = \sum_i r_i z^{i}$; it cannot recover any individual $a_i$ or
$r_i$. The assignment's statement ``I know coefficients corresponding to
$C_0,\dots,C_d$'' therefore fails the proof-of-knowledge requirement.
We fix this by running \emph{parallel} openings in a single combined
$\Sigma$-protocol, while still enforcing the evaluation constraint as a
scalar equation on the responses.

\section{The Parallel Per-Coefficient $\Sigma$-Protocol}

\subsection{Protocol}

\textbf{Statements verified:} (1) $g^{a_i} h^{r_i} = C_i$ for all $i$;
(2) $\sum_{i=0}^{d} a_i z^{i} \equiv y \pmod q$.

\begin{center}
\begin{tabular}{p{0.48\textwidth} c p{0.42\textwidth}}
\toprule
\textbf{Prover} & & \textbf{Verifier}\\
\midrule
For $i=0,\dots,d$: $\ k_{i,a}, k_{i,r} \xleftarrow{\$} \mathbb{Z}_q$ & & \\
$T_i \leftarrow g^{k_{i,a}} h^{k_{i,r}} \pmod p$ & & \\
$E \leftarrow \sum_{i=0}^{d} k_{i,a}\, z^{i} \pmod q$ & & \\
$\text{send } (T_0,\dots,T_d,\,E)$ & $\longrightarrow$ & \\
                                   & $\xleftarrow{c}$ & $c \xleftarrow{\$} \mathbb{Z}_q$ \\
For each $i$: & & \\
$\quad u_{i,a} = k_{i,a} + c\cdot a_i \pmod q$ & & \\
$\quad u_{i,r} = k_{i,r} + c\cdot r_i \pmod q$ & & \\
$\text{send } \{(u_{i,a},u_{i,r})\}_{i=0}^{d}$ & $\longrightarrow$ & \\
& & \textbf{Check (V1):} for all $i$,\\
& & $\quad g^{u_{i,a}} h^{u_{i,r}} \stackrel{?}{=} T_i \cdot C_i^{c} \pmod p$\\
& & \textbf{Check (V\_eval):}\\
& & $\quad \sum_{i} u_{i,a} z^{i} \stackrel{?}{\equiv} E + c\cdot y \pmod q$\\
& & Accept iff both checks pass. \\
\bottomrule
\end{tabular}
\end{center}

Repeat independently $k$ times with fresh nonces to amplify soundness.

\subsection{Why the evaluation check binds}

Computing $\sum_i u_{i,a}\, z^{i}$ honestly gives
\[
\sum_{i=0}^{d} u_{i,a}\, z^{i}
= \sum_{i=0}^{d}\bigl(k_{i,a}+c\,a_i\bigr)\, z^{i}
= \underbrace{\sum_{i} k_{i,a}\, z^{i}}_{=\,E}
+ c\,\underbrace{\sum_i a_i\, z^{i}}_{=\,P(z)}.
\]
Check (V\_eval) therefore passes iff $P(z)=y$. Crucially, $E$ is committed
in the first message, before $c$ is chosen, so the prover cannot adapt it.

\paragraph{Does $E$ leak?} $E=\sum_i k_{i,a}z^{i}$ is a fixed linear
combination of independent uniform random variables in $\mathbb{Z}_q$;
it is uniformly distributed on $\mathbb{Z}_q$, independent of the
coefficients $a_i$.

\section{Fiat--Shamir Non-Interactive Version}

Replace the verifier's random challenge with a hash of every public
element and the first message:
\[
c = H\bigl(\mathsf{tag} \,\|\, g \,\|\, h \,\|\, p \,\|\, q
\,\|\, C_0 \,\|\, \cdots \,\|\, C_d \,\|\, z \,\|\, y
\,\|\, T_0 \,\|\, \cdots \,\|\, T_d \,\|\, E\bigr) \bmod q,
\]
where $H=\mathrm{SHA\text{-}256}$, $\mathsf{tag}$ is a protocol
identifier, and each field is serialised as a $4$-byte big-endian length
prefix followed by its big-endian byte representation. Length prefixing
removes encoding ambiguity; the tag provides domain separation.

\paragraph{Security in the ROM.} Security follows from the Forking
Lemma of Pointcheval--Stern~\cite{ps2000}: a successful forger can be
rewound to produce two accepting transcripts sharing the same
$(T_0,\dots,T_d,E)$ but distinct challenges, from which the extractor
of Section~\ref{sec:security} recovers the full witness, contradicting
either binding or DL hardness.

\section{Security Analysis}\label{sec:security}

\begin{theorem}[Completeness]
An honest prover holding $(a_i,r_i)$ with $P(z)=y$ is always accepted.
\end{theorem}
\begin{proof}
\textbf{(V1):} $g^{u_{i,a}}h^{u_{i,r}} = g^{k_{i,a}+ca_i}h^{k_{i,r}+cr_i}
= (g^{k_{i,a}}h^{k_{i,r}})(g^{a_i}h^{r_i})^{c} = T_i\cdot C_i^{c}$.
\textbf{(V\_eval):} $\sum_i u_{i,a}z^{i} = \sum_i k_{i,a}z^{i}
+ c\sum_i a_i z^{i} = E + c\cdot P(z) = E + cy$.
\end{proof}

\begin{theorem}[Special Soundness --- Full Witness Extraction]
Given two accepting transcripts that share the same first message
$(T_0,\dots,T_d,E)$ but have distinct challenges $c\neq c'$:
\[
(T_\ast,E,c,\{(u_{i,a},u_{i,r})\})\quad\text{and}\quad
(T_\ast,E,c',\{(u'_{i,a},u'_{i,r})\}),
\]
one can efficiently extract the full witness
$(\hat a_0,\dots,\hat a_d,\hat r_0,\dots,\hat r_d)$ with
$C_i = g^{\hat a_i}h^{\hat r_i}$ for all $i$ and
$\sum_i \hat a_i z^{i} = y$.
\end{theorem}
\begin{proof}
Let $\delta = (c-c') \in \mathbb{Z}_q^{*}$. From (V1) on both transcripts:
$g^{u_{i,a}-u'_{i,a}} h^{u_{i,r}-u'_{i,r}} = C_i^{\,c-c'} = C_i^{\delta}$.
Setting
\[
\hat a_i = (u_{i,a}-u'_{i,a})\,\delta^{-1},\qquad
\hat r_i = (u_{i,r}-u'_{i,r})\,\delta^{-1}\quad (\bmod\ q)
\]
gives $g^{\hat a_i}h^{\hat r_i}=C_i$. Subtracting the two (V\_eval)
equations yields $\sum_i(u_{i,a}-u'_{i,a})z^{i} = (c-c')y = \delta y$;
dividing by $\delta$ gives $\sum_i \hat a_i z^{i} = y$. The extractor
outputs all $2(d+1)$ witness components.
\end{proof}

The soundness error per round is $1/q$, negligible in the bit-length of
$q$; $k$ independent rounds reduce it to $1/q^{k}$.

\begin{theorem}[Perfect Honest-Verifier Zero-Knowledge]
A PPT simulator without the witness produces transcripts identically
distributed to real transcripts.
\end{theorem}
\begin{proof}[Proof sketch]
Simulator: pick $c \xleftarrow{\$} \mathbb{Z}_q$; for each $i$ pick
$u_{i,a}, u_{i,r} \xleftarrow{\$} \mathbb{Z}_q$; define
$T_i \leftarrow g^{u_{i,a}}h^{u_{i,r}} C_i^{-c}$; define
$E \leftarrow \sum_i u_{i,a}z^{i} - cy$. Output
$(T_i, E, c, \{(u_{i,a},u_{i,r})\})$. Both (V1) and (V\_eval) hold by
construction. In a real run the responses $u_{i,a}=k_{i,a}+ca_i$ are
uniform and independent of $c$ because $k_{i,a}$ is; $T_i$ and $E$ are
uniquely determined by $(c,u_{i,a},u_{i,r},C_i,y)$ via the same formulae
as the simulator. Hence real and simulated distributions are identical.
\end{proof}

\section{Implementation}

All code is Python 3. The only standard-library imports are
\texttt{hashlib}, \texttt{secrets}, \texttt{random}, \texttt{json},
\texttt{os}, \texttt{sys}, and \texttt{time}. No external cryptographic
or ZK library is used, per the assignment constraint.

\subsection{File layout}

\begin{center}
\begin{tabular}{ll}
\toprule
\textbf{File} & \textbf{Role} \\
\midrule
\texttt{generate\_inputs.py}   & Input generator (primes, generators, witness)\\
\texttt{Problem4.py}           & NI and interactive-hash prover/verifier\\
\texttt{prover\_interactive.py}    & Live two-program prover (Variant B)\\
\texttt{verifier\_interactive.py}  & Live two-program verifier (Variant B)\\
\texttt{public.json}           & Public inputs (read by both sides)\\
\texttt{private.json}          & Prover's witness\\
\texttt{proof.json}            & Proof transcript (prover $\to$ verifier)\\
\bottomrule
\end{tabular}
\end{center}

\subsection{Core $\Sigma$-protocol}

\begin{lstlisting}[style=pythonstyle, caption={Commit, response, and verification.}]
def commit_phase(d, z, g, h, p, q):
    k_a = [secrets.randbelow(q) for _ in range(d + 1)]
    k_r = [secrets.randbelow(q) for _ in range(d + 1)]
    T_list = [(pow(g, ka, p) * pow(h, kr, p)) % p
              for ka, kr in zip(k_a, k_r)]
    E, zp = 0, 1
    for ka in k_a:
        E  = (E + ka * zp) % q
        zp = (zp * z)      % q
    return T_list, E, k_a, k_r

def response_phase(c, coeffs, randomness, k_a, k_r, q):
    u_a = [(ka + c * a) % q for ka, a in zip(k_a, coeffs)]
    u_r = [(kr + c * r) % q for kr, r in zip(k_r, randomness)]
    return u_a, u_r

def verify_checks(C_list, z, y, T_list, E, c, u_a, u_r, g, h, p, q):
    for C, T, ua, ur in zip(C_list, T_list, u_a, u_r):
        lhs = (pow(g, ua, p) * pow(h, ur, p)) % p
        rhs = (T * pow(C, c, p)) % p
        if lhs != rhs:
            return False
    lhs_eval, zp = 0, 1
    for ua in u_a:
        lhs_eval = (lhs_eval + ua * zp) % q
        zp       = (zp * z)             % q
    return lhs_eval == (E + c * y) % q
\end{lstlisting}

\subsection{Fiat--Shamir challenge}

\begin{lstlisting}[style=pythonstyle, caption={Canonical length-prefixed SHA-256 challenge.}]
def _enc_int(n):
    raw = n.to_bytes((n.bit_length() + 7) // 8 or 1, "big")
    return len(raw).to_bytes(4, "big") + raw

def fs_challenge_ni(g, h, p, q, commitments, z, y, T_list, E):
    H = hashlib.sha256()
    H.update(b"ZKP-POLY-FS-v1:")
    for v in (g, h, p, q):
        H.update(_enc_int(v))
    H.update(_enc_int(len(commitments)))
    for C in commitments:
        H.update(_enc_int(C))
    H.update(_enc_int(z))
    H.update(_enc_int(y))
    for T in T_list:
        H.update(_enc_int(T))
    H.update(_enc_int(E))
    return int.from_bytes(H.digest(), "big") % q
\end{lstlisting}

\subsection{Two interactive variants}

\paragraph{Variant A --- hash-derived per-round challenge.} A single
program runs $k$ rounds, using a separate domain tag and round index
in the hash:
\[
c_j = H\bigl(\mathsf{"ZKP\text{-}POLY\text{-}INT\text{-}v1:"} \,\|\, j
\,\|\, \text{public} \,\|\, T_0^{(j)} \,\|\, \dots \,\|\, T_d^{(j)}
\,\|\, E^{(j)}\bigr) \bmod q.
\]
The verifier recomputes $c_j$ independently, so a malicious prover
cannot grind it. This matches the soundness of the Fiat--Shamir variant
per round.

\paragraph{Variant B --- live two-program exchange.} The prover and
verifier are separate processes communicating via a file-based message
queue under \texttt{./session/}. Per round $j$ the prover writes
\texttt{round\_j\_A.json} (commit), the verifier samples
$c_j \xleftarrow{\$} \mathbb{Z}_q$ and writes \texttt{round\_j\_B.json}
(challenge), the prover reads it and writes \texttt{round\_j\_C.json}
(response), and the verifier writes \texttt{round\_j\_V.json} (per-round
verdict). This faithfully realises the classical $\Sigma$-protocol
message flow.

\section{Input Generator}

\texttt{generate\_inputs.py} produces \texttt{public.json} and
\texttt{private.json} entirely from scratch. Allowed imports are
\texttt{hashlib} (SHA-256 only), \texttt{secrets}, \texttt{random},
\texttt{json}, and \texttt{sys}; all cryptographic primitives are
hand-written.

\subsection{Miller--Rabin primality test}

\begin{lstlisting}[style=pythonstyle]
def is_prime(n):
    if n < 2: return False
    for p in _SMALL_PRIMES:
        if n == p: return True
        if n % p == 0: return False
    d, r = n - 1, 0
    while d % 2 == 0:
        d //= 2; r += 1
    bases = _DETERMINISTIC_BASES if n < (1 << 82) \
            else [2 + secrets.randbelow(n - 3) for _ in range(40)]
    for a in bases:
        if a % n == 0: continue
        x = pow(a, d, n)
        if x == 1 or x == n - 1: continue
        composite = True
        for _ in range(r - 1):
            x = x * x % n
            if x == n - 1: composite = False; break
        if composite: return False
    return True
\end{lstlisting}

Using the deterministic base set $\{2,3,5,7,11,13,17,19,23,29,31,37\}$
the test is \emph{exact} for all $n < 3.3\cdot10^{24}$; above that we
fall back to $40$ random bases (error probability $\le 2^{-80}$).

\subsection{Safe-prime search and generators}

We sample a random $b$-bit odd integer $q$, prefilter by small-prime
trial division on both $q$ and $p = 2q+1$, then Miller--Rabin-test each.
A generator of the order-$q$ subgroup is obtained as $g = x^{2} \bmod p$
for random $x$, verifying $g \ne 1$ and $g^{q} \equiv 1$.

The second generator $h$ is produced by one of two user-selected methods:

\begin{itemize}
  \item[\texttt{--h-mode hash}]
    Hash-to-subgroup: $h = x^{2} \bmod p$ where
    $x = \mathrm{SHA256}(\mathsf{"ZKP\text{-}POLY\text{-}H\text{-}v1"} \| \text{ctr})
    \bmod p$, with counter incremented until the result has order $q$
    and is distinct from $\{0,1,g\}$. No one (prover, verifier, or
    designer) knows $\log_g h$.
  \item[\texttt{--h-mode discard}]
    Sample $s \xleftarrow{\$} \mathbb{Z}_q^{*}$, set $h = g^{s} \bmod p$,
    then discard $s$. This is trapdoor in principle (whoever ran the
    generator once knew $s$), so the hash-to-subgroup method is
    preferred in production.
\end{itemize}

\subsection{Parameter-size flag}

\begin{center}
\begin{tabular}{lll}
\toprule
\textbf{Flag} & \textbf{Size of $q$} & \textbf{Purpose} \\
\midrule
\texttt{--size toy}   & $q=101$ (PDF's example)       & Human-readable grading \\
\texttt{--size small} & $128$ bits                    & Fast end-to-end runs \\
\texttt{--size full}  & $256$ bits                    & Production demo \\
\bottomrule
\end{tabular}
\end{center}

\section{Testing and Results}

\subsection{End-to-end test matrix}

All six $\{\text{toy, small, full}\}\times\{\text{hash, discard}\}$
combinations produce valid inputs and pass every proof mode. A tampered
proof (where the first response scalar is shifted by $+1 \bmod q$) is
correctly rejected in every case.

\begin{center}
\begin{tabular}{llccc}
\toprule
\textbf{size} & \textbf{h-mode} & \textbf{NI} & \textbf{interactive-hash} & \textbf{NI-tampered} \\
\midrule
toy   & hash    & ACCEPT & ACCEPT & REJECT \\
toy   & discard & ACCEPT & ACCEPT & REJECT \\
small & hash    & ACCEPT & ACCEPT & REJECT \\
small & discard & ACCEPT & ACCEPT & REJECT \\
full  & hash    & ACCEPT & ACCEPT & REJECT \\
full  & discard & ACCEPT & ACCEPT & REJECT \\
\bottomrule
\end{tabular}
\end{center}

The live two-program variant (\texttt{prover\_interactive.py} /
\texttt{verifier\_interactive.py}) was also exercised with $k=4$ rounds
on 128-bit parameters and produced ACCEPT.

\subsection{Concrete toy example}

\begin{lstlisting}[style=jsonstyle, caption={\texttt{public.json} (toy)}]
{
  "p": 607, "q": 101, "g": 7, "h": 49,
  "d": 2,   "k": 3,
  "commitments": [346, 270, 552],
  "z": 9,   "y": 5
}
\end{lstlisting}

\begin{lstlisting}[style=jsonstyle, caption={\texttt{private.json} (toy)}]
{ "coeffs":     [33, 92, 63],
  "randomness": [70, 39, 10] }
\end{lstlisting}

Verification: $P(9) = 33 + 92\cdot 9 + 63\cdot 81 = 33 + 828 + 5103 = 5964
\equiv 5 \pmod{101}$. \checkmark

\begin{lstlisting}[style=jsonstyle, caption={\texttt{proof.json}
(non-interactive, after running \texttt{Problem4.py prover --mode noninteractive})}]
{
  "mode": "noninteractive",
  "transcript": {
    "T_list": [552, 212, 26],
    "E":      9,
    "c":      55,
    "u_a":    [18, 54, 11],
    "u_r":    [43, 36, 22]
  },
  "tampered": false
}
\end{lstlisting}

\textbf{Verifier output:} \texttt{ACCEPT}. Manually check (V\_eval):
$\sum_i u_{i,a}\cdot 9^{i} = 18 + 54\cdot 9 + 11\cdot 81
 = 18 + 486 + 891 = 1395 \equiv 82 \pmod{101}$, and
$E + cy = 9 + 55\cdot 5 = 284 \equiv 82 \pmod{101}$. \checkmark

\subsection{REJECT example}

Re-running the prover with \texttt{--tamper} shifts $u_{0,a}$ by $+1
\bmod q$. The verifier prints \texttt{REJECT}.

\subsection{Reproduction}

\begin{lstlisting}[style=bashstyle]
# 1. generate inputs
python generate_inputs.py --size small --h-mode hash \
                          --degree 3 --rounds 4

# 2. non-interactive proof
python Problem4.py prover   --mode noninteractive
python Problem4.py verifier                       # ACCEPT

# 3. single-program interactive (hash-derived challenge)
python Problem4.py prover   --mode interactive-hash
python Problem4.py verifier                       # ACCEPT

# 4. live two-program interactive
python verifier_interactive.py &                  # shell A
python prover_interactive.py                      # shell B

# 5. tamper demo
python Problem4.py prover   --mode noninteractive --tamper
python Problem4.py verifier                       # REJECT
\end{lstlisting}

\section{Security Summary}

\begin{center}
\begin{tabular}{p{0.27\textwidth} p{0.67\textwidth}}
\toprule
\textbf{Property} & \textbf{Statement}\\
\midrule
Completeness      & Honest $(a_i,r_i)$ with $P(z)=y$ always pass (V1) and (V\_eval).\\
Special Soundness & Two transcripts with $(c\ne c')$ on same first message extract the \emph{full} $2(d+1)$-dim witness.\\
Perfect HVZK      & Real and simulated transcripts are identically distributed.\\
Evaluation Binding& (V\_eval) forces $\sum_i a_i z^{i}=y$; cheating requires either a false $E$ (caught by extraction) or breaking DL.\\
Non-interactive   & Fiat--Shamir in ROM yields a NIZK of equivalent security via the Forking Lemma.\\
\bottomrule
\end{tabular}
\end{center}

\section{Conclusion}

We delivered a from-scratch implementation of a zero-knowledge proof of
knowledge and correct evaluation of a secret polynomial, satisfying all
requirements of the assignment: Pedersen commitments, a multi-round
interactive $\Sigma$-protocol, a Fiat--Shamir non-interactive variant,
rigorous completeness / special-soundness / HVZK proofs, and an input
generator built without external cryptographic libraries. The design
extracts the \emph{full} witness (not just an aggregated linear
combination), which strictly meets the proof-of-knowledge standard.

\section*{References}
\begin{enumerate}
  \item T.~P.~Pedersen, ``Non-Interactive and Information-Theoretic
        Secure Verifiable Secret Sharing,'' \textit{CRYPTO 1991}.
  \item C.~P.~Schnorr, ``Efficient Identification and Signatures for
        Smart Cards,'' \textit{CRYPTO 1989}.
  \item M.~Bellare and P.~Rogaway, ``Random Oracles are Practical,''
        \textit{ACM CCS 1993}.
  \item \label{ps2000} D.~Pointcheval and J.~Stern, ``Security Proofs
        for Signature Schemes,'' \textit{Journal of Cryptology}, 2000.
  \item D.~Boneh and V.~Shoup, \textit{A Graduate Course in Applied
        Cryptography}, 2023.
\end{enumerate}

\section*{Anti-Plagiarism Statement}
We certify that this submission is entirely our own work. All references
are cited appropriately. No external cryptographic library was used.

\vspace{1cm}
\textbf{Names:} Divyansh Sevta, Nishanth D. \qquad \textbf{Date:} \today
\end{document}
